% Adapted from Rover Resume - Base Template
% Link: https://github.com/subidit/rover-resume
% License: CC BY 4.0

\documentclass[10pt]{article}
\usepackage[margin=0.65in, a4paper]{geometry}
\usepackage[bookmarks=false]{hyperref}
\hypersetup{
  colorlinks=true,
  urlcolor=black,
  pdfauthor={Seoung Choi},
  pdftitle={Seoung Choi CV}
}
\setcounter{secnumdepth}{0}
\pagestyle{empty}
\pdfgentounicode=1

\newcommand{\cvsection}[1]{%
  \vspace{0.55em}
  \noindent{\large\bfseries\MakeUppercase{#1}}\par
  \vspace{0.15em}
  \hrule
  \vspace{0.35em}
}

\newcommand{\cventry}[4]{%
  \noindent\textbf{#1}\hfill #2\par
  \noindent\textit{#3}\hfill #4\par
}

\begin{document}

\begin{center}
\begin{minipage}{0.52\textwidth}
{\Huge\bfseries Seoung Choi}\\
\medskip
M.S. Student, School of Computing, KAIST
\end{minipage}
\hfill
\begin{minipage}{0.42\textwidth}
\raggedleft
Email: \href{mailto:choisw0823@kaist.ac.kr}{choisw0823@kaist.ac.kr}\\
Website: \href{https://choisw0823.github.io}{choisw0823.github.io}\\
GitHub: \href{https://github.com/choisw0823}{github.com/choisw0823}
\end{minipage}
\end{center}

\cvsection{Education}
\noindent\textbf{Korea Advanced Institute of Science \& Technology (KAIST), School of Computing}\par
\noindent\textit{M.S. Student}\hfill Sep. 2026 -- Present\par
\vspace{0.2em}
\noindent\textbf{Korea Advanced Institute of Science \& Technology (KAIST), School of Computing}\par
\noindent\textit{B.S.}\hfill Mar. 2021 -- Feb. 2026\par

\cvsection{Research Interests}
\begin{itemize}
  \item AI agents for automating complex multi-step tasks, video understanding, and multimodal AI systems that better understand, reason over, and interact with visual information.
\end{itemize}

\cvsection{Experience}
\cventry{MLV Lab, KAIST}{Daejeon, Korea}{Intern}{Jun. 2025 -- Present}
\begin{itemize}
  \item Conduct research in multimodal AI and vision-language models, advised by Prof. Hyunwoo J. Kim.
\end{itemize}

\cventry{Smoretalk}{Korea}{AI Engineer}{Dec. 2024 -- Aug. 2025}
\begin{itemize}
  \item Participated in AI service development using image generation and editing technologies.
\end{itemize}

\cventry{Naviworks, DX Technology Team}{Korea}{Internship, AI Engineer}{Feb. 2024 -- Aug. 2024}
\begin{itemize}
  \item Participated in the metaieyes digital twin platform project, developing product features for converting 2D data into 3D data.
\end{itemize}

\cvsection{Publications}
\begin{enumerate}
  \item Seohyun Lee*, \textbf{\underline{Seoung Choi}}*, Dohwan Ko*, Jongha Kim, and Hyunwoo J. Kim. ``VideoSearch-R1: Iterative Video Retrieval and Reasoning via Soft Query Refinement.'' ECCV, 2026. (*Equal contribution)
  \item Dohwan Ko, Jinyoung Park, \textbf{\underline{Seoung Choi}}, Sanghyeok Lee, Seohyun Lee, and Hyunwoo J. Kim. ``MoE-GRPO: Optimizing Mixture-of-Experts via Reinforcement Learning in Vision-Language Models.'' CVPR, 2026.
\end{enumerate}

\cvsection{Patent}
\begin{itemize}
  \item Apparatus and method for generating standard objects based on satellite imagery. Korean Patent Application No. 10-2024-0117862, filed Aug. 30, 2024; Publication No. 10-2026-0031996, published Mar. 9, 2026.
\end{itemize}

\cvsection{Skills}
\begin{description}
  \item[Research] AI agents, video understanding, multimodal AI, vision-language models, retrieval and reasoning.
  \item[Engineering] Python, PyTorch, Git, model training and evaluation, AI service prototyping.
\end{description}

\end{document}
